%% kapa: exam=1 round=20 year=2009 period=1 form=A subject=civil_law subject_ko=민법 seq=31 number=31
%% source: Q-Net 감정평가사 기출문제 (cst00309) artlSeq=5024178 file=extracted/1008/5024178/2009년 제20회 감정평가사 1차시험 기출문제/1교시.hwp
%% conversion: kordoc 4.24.51 -> tools/gen1.py (auto); review=proofread-sonnet; flags=none
\begin{question}{31}
다음 설명 중 옳지 \uline{않은} 것은?
\begin{choices}
  \item 조건성취의 효과는 원칙적으로 소급하지 않는다.
  \item 기간의 말일이 토요일 또는 공휴일인 때에는 기간은 그 익일로 만료한다.
  \item 시기부 법률행위는 기한이 도래한 때부터 효력이 생기고 종기부 법률행위는 기한이 도래한 때부터 그 효력을 잃는다.
  \item 기한의 이익은 특별한 사정이 없는 한, 채무자의 이익을 위한 것으로 추정한다.
  \item 기간의 계산에 있어 연령의 계산에는 출생일을 산입하지 않는다.
\end{choices}
\end{question}
\answer{5}
